\section{Synthesis and Proof Boundaries}
\label{sec:synthesis}

The checked results provide several components of a possible persistence-based dynamics: finite survival selection, phase-covariant transport, dispersion/Clifford equivalence, geometric gauge isotropy, conditional unitary reconstruction, scalar sign selection and a candidate two-sector fermion mechanism. They are not yet a derivation of a unique complete physical action.

\subsection{Leading Operators versus Exact Exhaustion}
Within an assumed dimension-four list, \nolinkurl{Operators.allowed_monomials} bounds the number of fields and derivatives. This classification does not prove the physical cutoff itself, determine representations, or exclude all other invariants. The actual cosine action and the finite-amplitude scalar counterexample demonstrate why higher terms cannot be deleted solely by invoking finite capacity.

The controlled $D_4$ gauge limit is a positive replacement for exact truncation: the leading Abelian constant-field response follows with a bound on each phase's error. Nonconstant and non-Abelian response, fermion continuum control and gravitational dynamics need analogous proofs.

\subsection{Matter and Radiation}
\label{sec:matter_radiation}
The product $\xi\mathbf K f$ motivates a complexity/rate trade-off under a fixed budget. By itself it does not classify physical excitations into exactly two stable phases. A classification of massive and massless modes requires the actual update's spectrum, conserved observables and stability conditions. The candidate domain-wall construction is one step toward such a spectrum; it does not prove the proposed universal matter/radiation bifurcation.

\subsection{Connection to Numerical Constants}
A field normalization, scale convention and common response must be fixed before operator coefficients are identified with measured constants. Paper III tests its proposed formulas jointly. Improving a storage-overhead mechanism cannot substitute for closing the scalar, fermion and gauge response relations at a common normalization.
